\section{LaMET 与混合重正化的表述}
\label{sec:analytic}

强子的领头扭度（leading-twist）胶子 PDF $g(x)$ 定义为~\cite{Collins:1989gx}
%
\begin{align}\label{eq:LCPDF}
g(x,\mu) = \frac{1}{ 2 x P^+} \int_{-\infty}^\infty \frac{\mathrm{d} \xi^-}{2\pi} \,& e^{-i x P^+ \xi^-} \langle P | F_a^{+\mu}(\xi^-) \notag\\
&
 \times \mathcal{W}_{ab}(\xi^-, 0) \, F_{b,\mu}^+(0) | P \rangle,
\end{align}
其中 $|P\rangle$ 表示沿 $z$ 方向具有动量 $P$ 的非极化强子。这里，$x$ 是胶子携带的纵向动量分数，$\mu$ 表示 $\overline{\rm MS}$ 方案中的重正化标度。光锥坐标由 $\xi^\pm=(\xi^t\pm \xi^z)/\sqrt{2}$ 给出，规范不变性由 Wilson 线保证：
\begin{align}
\mathcal{W}_{ab}(\xi^-, 0)= \mathcal{P}\,{\rm exp}\big[ -i g \int_0^{\xi^-} \mathrm{d} \eta^- \, A_c^+(\eta^-) t^c \big]_{ab},
\end{align}
其中指标 $\{a,b\}$ 表示 SU(3) 伴随表示的色分量，而 $t^c$ 是规范群的生成元。

按照 LaMET，胶子 PDF 可以通过计算一个空间关联的矩阵元来提取：
\begin{align}
h^{\text{B}}(z,P_z)= \langle P | &F^{ti}(z)\mathcal{W}(z,0)F^t_{i}(0) \notag\\
&-F^{ij}(z)\mathcal{W}(z,0)F_{ij}(0)| P \rangle,
\end{align}
其中 $i$、$j$ 仅遍历横向指标 $\{x,y\}$。
下文为记号简洁起见，我们约定省略色指标。

尽管定域情形并不严格容许乘法重正化，我们并未纳入来自定域区域的信息，因为结果是归一化的，因而对定域贡献不敏感。非定域的裸准光前（light-front, LF）关联可以乘法重正化为~\cite{Zhang:2018diq}
%
\begin{align}
h^{\text{B}}(z,P_z)=Z_L e^{-\bar{\delta} m |z|} \,h^{\text{R}}(z,P_z).
\end{align}
其中 $h^{\text{R}}(z,P_z)$ 是重正化的准光前关联，它同时包含由 $Z_L$ 刻画的与 $z$ 无关的对数紫外（UV）发散，以及来自 Wilson 线自能的、随 $z$ 变化的线性发散，后者记作 $\bar{\delta} m$。
为了消除这些发散，我们采用混合重正化方案~\cite{Ji:2020brr}，将短距离与长距离分别处理：
在短距离处，我们对静止系矩阵元施加比值方案（ratio scheme）~\cite{Radyushkin:2018cvn}；
而在长距离处，我们通过拟合零动量矩阵元 $h^{\text{B}}(z,0)$ 进行质量重正化，以捕捉线性发散效应。
所得的重正化准光前关联表示为：
\begin{equation}\label{eq:hybridscheme}
h^{\mathrm{R}}(z, P_z) =
\begin{cases}
\!\displaystyle\!
\frac{h^{\mathrm{B}}(0,0)}{h^{\mathrm{B}}(0,P_z)} \frac{h^{\mathrm{B}}(z,P_z)}{h^{\mathrm{B}}(z,0)} & z \leq z_s \\[12pt]
\!\displaystyle\!
\frac{h^{\mathrm{B}}(0,0)}{h^{\mathrm{B}}(0,P_z)}  \frac{h^{\mathrm{B}}(z,P_z)}{h^{\mathrm{B}}(z_s,0)} e^{(\delta m + m_0)(z - z_s)} & z > z_s
\end{cases}
\end{equation}
其中 $z_s$ 区分短距离与长距离，$\delta m$ 与 $m_0$ 分别对应质量重正化和 renormalon（重整子）模糊性，引入它们是为了抵消来自 Wilson 线自能的线性发散。在 $z_s \to \infty$ 的极限下，该表达式退化为忽略 Wilson 线自能贡献的标准比值方案重正化。
\UPDATE{在短距离采用双比值是处理 $z \rightarrow 0$ 极限与格距 $a \rightarrow 0$ 极限不可交换性所导致的离散化误差的好方法~\cite{Ji:2020brr}，不过这一方案会使所得胶子 PDF 由其第一矩 $\langle x\rangle_g$ 归一化。
在先前关于胶子 PDF 连续极限的研究中可以看到，用比值方案处理短距离时离散化误差在当前统计误差内控制得很好~\cite{Fan:2022kcb}。
有建议指出，采用自重正化（self-renormalization）方法可以更好地控制离散化效应~\cite{LatticePartonLPC:2021gpi}，但这需要使用多个格距的数据。}

\begin{figure*}[htbp]
\centering
        \includegraphics[width=.48\textwidth]{images/W3_m0vz.pdf}
\centering
        \includegraphics[width=.48\textwidth]{images/H5_m0vz.pdf}

    \caption{$\mu = 2 $ GeV 处各强子在两种涂抹方式（左：Wilson-3；右：HYP5）下，$\delta m + m_0$ 随中心点 $z$ 变化的拟合结果。
    误差带仅反映统计涨落。
    各图中的数据点代表用于后续混合重正化的 $\delta m + m_0$ 取值。}
    \label{fig:m0_vs_z}
\end{figure*}

特别地，若拥有多个格距，便可通过拟合分别确定 $\delta m$ 和 $m_0$~\cite{Ji:2020brr, LatticePartonLPC:2021gpi, Zhang:2023bxs}。
然而在单一格距的情形下，更实用的做法是将 $P_z = 0$ 的裸矩阵元与微扰计算的 ``Wilson 系数" 相匹配，只拟合二者之和 $\delta m + m_0$。该 Wilson 系数可由文献~\cite{Balitsky:2021qsr} 导出，为
\begin{align}\label{eq:wil-coef}
\mathcal{H}(z,\mu) = 1 + \frac{\alpha_s C_A}{2\pi} \left( \frac{5}{6} \ln \left( \frac{z^2 \mu^2 e^{2\gamma_E}}{4} \right) + \frac{3}{2} \right).
\end{align}
随后我们在短距离区域通过拟合如下关系来确定组合参量 $\delta m + m_0$：
\begin{equation}\label{eq:dm_fit}
(\delta m + m_0) z  + I_0 \approx \ln\left[ \frac{\mathcal{H}(z,\mu)}{h^{\mathrm{B}}(z,0)} \right],
\end{equation}
其中 $I_0$ 表示一个与 $z$ 无关的常数。

重正化的准光前关联函数可经傅里叶变换转化为所谓的准 PDF。胶子准 PDF 通过如下的因子化公式与光前 PDF 相联系：
\begin{align}\label{eq:mtcheq}
 x\tilde{g}\left(x, P_z\right)= & \int_{-1}^1 \frac{d y}{|y|} C\left(\frac{x}{y}, \frac{\mu}{y P_z}\right)  y g(y, \mu)  +h.t. ,
\end{align}
其中 $g(y, \mu)$ 是重正化标度 $\mu$ 处的光前胶子 PDF，$C\left(x/y, \mu/(y P_z)\right)$ 表示微扰匹配系数，而 $h.t.$ 表示按 $1/P_z$ 幂次压低的高扭度修正。
\UPDATE{注意，本工作忽略了与夸克混合的贡献，因为赝 PDF 研究~\cite{Fan:2022kcb}已表明其影响小于约 $10\%$，远在统计精度之内，也小于本文所探讨的格点系统效应。}
比值方案下坐标空间中至 NLO 的微扰匹配系数可在文献~\cite{Balitsky:2019krf, Balitsky:2021qsr} 中找到。随后我们按照文献~\cite{Yao:2022vtp} 所述步骤，直接经傅里叶变换得到对应的动量空间匹配核。比值方案与混合方案的匹配系数给出如下：

\begin{equation}\label{eq:kernel}
C^{\text{ratio}}\left(\xi, \frac{\mu}{p_z}\right)
= \delta(1-\xi)
+ \frac{\alpha_s C_A}{2\pi}
\left\{
\begin{array}{ll}
\left(
2 \frac{(1 - \xi + \xi^2)^2}{1 - \xi} \ln\left(\frac{\xi}{\xi-1}\right)
+\frac{11 - 28\xi + 18\xi^2 - 12\xi^3}{6(1-\xi)}
\right)_{+(1)}, & \xi > 1 \\[1.5ex]
\left(
2 \frac{(1 - \xi + \xi^2)^2}{1 - \xi} \left[ -\ln(\frac{\mu^2}{4{p_z}^2}) + \ln\left[\xi(1-\xi)\right] \right]
- \frac{15 - 56\xi + 102\xi^2 - 96\xi^3 + 48\xi^4}{6(1-\xi)}
 \right)_{+(1)}, & 0 < \xi < 1 \\[1.5ex]
\left(
-2 \frac{(1 - \xi + \xi^2)^2}{1 - \xi} \ln\left(\frac{-\xi}{1-\xi}\right)
- \frac{11 - 28\xi + 18\xi^2 - 12\xi^3}{6(1-\xi)} \right)_{+(1)}, & \xi < 0
\end{array}
\right.
\end{equation}

\begin{equation}
C^{\text{hyb.}}\left(\xi, \frac{\mu}{p_z}\right) =C^{\text{ratio}}\left(\xi, \frac{\mu}{p_z}\right) + \frac{\alpha_s C_A}{2\pi} \frac 5 6 \bigg[-\frac{1}{|1-\xi|} +\frac{2\text{Si}((1-\xi)z_s p_z)}{\pi(1-\xi)}\bigg]_+ .
\end{equation}

\noindent 其中 $\xi = \frac{x}{y}$，而 $p_z$ 表示部分子携带的动量，即 $y P_z$。
